\textbf{General instructions.} Scientists have discovered a crashed alien
flying saucer. High up inside the hold, they found several screens
transmitting views of the sky and telescopes have been set up to let you see
them clearly from the level of the deck. Use your telescope to observe the
(simulated) targets on the screens and record your results.

There are 5 screens on the opposite side: the central one will display video
for tasks 1 and 2, the other four will display static images for tasks 3 and
4. The two screens closer to the centre will display the (same) image for
task 3, and the two outer screens will display the (same) image for task 4.
Point your telescope at the screens furthest away from you.

You will have a total of 30 minutes to complete the observing tasks, however
tasks 1 and 2 will only be displayed once: just as with real observations
you will only have one opportunity to collect the data. There will be two
clocks visible showing the time remaining in the round.

At the start of the round a clock on the central screen will show the
simulated time at the observer's location. The clock will have the correct
orientation when seen through the telescope. The time will be shown for 3
minutes after which it will disappear; use this to set a start time for your
observations.

Caution: the scale of the field of view is different between the video and
still images.

\medskip
\textbf{Observation 1: `Asteroid occultation'} (15 points)

Calculations based on the orbital elements predict that an asteroid will
occult the star HD 163390 for 21 s, with the maximum occultation (mid-time)
occurring at 23:03:32 UT. However, the ephemeris is not perfect and the
prediction may be wrong by up to 20 s for the time and by 10 s for the
duration.

Based on your observations, find the true mid-time and duration of the
occultation. To identify the star use Map 1 and the following coordinates:
\[ \mathrm{HD\ 163390:}\quad \mathrm{RA} = 17^h58^m05^s, \quad
   \mathrm{DEC} = -18^\circ 50' 46.14'' \]
The map and the sky are in the same epoch.

[Answer table: mid-time of occultation $\pm$ error; duration of occultation
$\pm$ error]

% FIGURE NEEDED: Map no. 1 -- star chart of the field around HD 163390

\medskip
\textbf{Observation 2: `Starlink'} (15 points)

In the same star field as for Question 1, a `train' of Starlink satellites
will appear near the meridian of $17^h59^m$ at around 23:05 UT. Their
passage will last for around three minutes.

You may assume that the centre of the star field is at an altitude of
$20^\circ$ and that the satellites are 400 km above the Earth's surface
moving on circular orbits with equal distances between them. You may also
assume that satellites will move vertically (perpendicular to the horizon).

\begin{description}
\item[(a)] Measure the angular velocity of the satellites as seen by an
  observer on the simulated sky.
\item[(b)] Measure the time interval between the passes of successive
  satellites and mark their path on the sky chart (Map 1).
\item[(c)] Calculate the theoretical angular velocity of the satellites as
  seen by the observer, using the information given in the question.
\item[(d)] Estimate the distance in km between two consecutive satellites.
\end{description}

Constants: $G = 6.674 \times 10^{-11}\,\mathrm{N\,m^2\,kg^{-2}}$;
$M_\oplus = 5.972 \times 10^{24}$\,kg; $R_\oplus = 6378$\,km.

\medskip
\textbf{Observation 3: `Planetary Moons'} (10 points)

The screen will display an image of one of the planets of the Solar System
as seen on August 15, 2023, at 00:00 UT. Identify any five moons and mark
them on the answer sheet (you may use the moon position chart attached below
and the table showing their brightness).

% FIGURE NEEDED: moon position chart (the numbers on the left indicate the
% days of August 2023 at midnight UT; moon numbers I, II, ... as in the
% table below)

\begin{center}
\begin{tabular}{cll}
Number & Name & Magnitude \\
\hline
I & Mimas & 13.0 \\
II & Enceladus & 11.8 \\
III & Tethys & 10.4 \\
IV & Dione & 10.6 \\
V & Rhea & 9.9 \\
VI & Titan & 8.5 \\
VII & Hyperion & 14.4 \\
VIII & Japetus & 11.0 \\
\end{tabular}
\end{center}

[Answer sheet: image of the planet on which the positions of any 5 moons are
to be marked with a dot and labelled with their numbers (I, II, \ldots)]

\medskip
\textbf{Observation 4: `Supernova'} (10 points)

The other screen presents the view of a galaxy and a bright (mag $< 11$)
object which was not visible previously. Estimate the right ascension (RA)
and declination (DEC) coordinates of this star and estimate its magnitude.
You may use Map 2, with stellar coordinates and a list of magnitudes.

[Answer table: right ascension, declination, estimated magnitude]

% FIGURE NEEDED: Map no. 2 -- star chart of the galaxy field with stellar
% coordinates and a list of magnitudes
